\chapter{Explicabilité et rapport de recommandations}\label{chapExpl}
\pagestyle{fancy}
\lhead{Explicabilité et rapport de recommandations}

\graphicspath{{Chapitre4/images_chapitre4/}{images_chapitre4/}{Chapitre3/images_chapitre3/}{images_chapitre3/}}

% Couleurs de tableaux (déjà définies dans le préambule principal : ne sont créées que si absentes)
\providecolor{thead}{HTML}{1F3A5F}
\providecolor{trowa}{HTML}{EEF2F7}

% Insère une figure si le fichier existe, sinon un cadre « à insérer »
\providecommand{\figtrois}[2]{%
  \IfFileExists{Chapitre3/images_chapitre3/#1}{\includegraphics[width=#2\textwidth]{#1}}{%
  \IfFileExists{images_chapitre3/#1}{\includegraphics[width=#2\textwidth]{#1}}{%
  \IfFileExists{Chapitre4/images_chapitre4/#1}{\includegraphics[width=#2\textwidth]{#1}}{%
  \IfFileExists{images_chapitre4/#1}{\includegraphics[width=#2\textwidth]{#1}}{%
  \fbox{\parbox[c][4cm][c]{#2\textwidth}{\centering\small\itshape
  Figure à insérer : \texttt{\detokenize{#1}}}}}}}}}

% =====================================================================
\section{Introduction}

Un modèle qui répond seulement « attaque » ou « normal » ne suffit pas à un
opérateur de chaîne du froid : avant de bloquer un lot de médicaments, il
doit savoir \emph{pourquoi} le modèle soupçonne une falsification, puis
\emph{quoi faire}. Ce chapitre apporte ces deux réponses, dans l'ordre où
l'application les produit :

\begin{itemize}
  \item l'\textbf{explicabilité} des décisions du modèle retenu, par la
        méthode SHAP (section~\ref{sec:expl_shap}) : quelles variables font
        basculer le score vers « attaque », en général et pour chaque
        scénario ;
  \item le \textbf{rapport de recommandations} destiné à l'opérateur
        (section~\ref{sec:expl_rapport}) : pour chaque scénario d'attaque,
        son niveau de détection mesuré au test, les indices utilisés par le
        modèle et les consignes à appliquer.
\end{itemize}

Les deux parties sont enchaînées : le classement des variables donné par
SHAP est l'une des entrées du rapport, qui est servi par l'API et affiché par
la page « Rapport pour l'opérateur » du tableau de bord
(chapitre~\ref{chap5}).

Un point de méthode conditionne tout le chapitre : le modèle est
\textbf{binaire}. Il est entraîné sur l'étiquette « normal » ou « attaque » et
ne nomme jamais l'attaque. Les neuf scénarios cités dans le rapport sont ceux
que nous avons \emph{simulés}, donc connus au moment de l'évaluation : ils
servent à mesurer comment le modèle réagit à chacun. En service, le type
d'attaque est inconnu ; la famille d'attaque est alors \emph{déduite} des
indices qui ont le plus pesé dans le score (section~\ref{sec:expl_deduction}),
sous forme d'hypothèse à confirmer sur place.


% =====================================================================
\section{Explicabilité du modèle par SHAP}
\label{sec:expl_shap}

\subsection{Principe des valeurs de Shapley}

SHAP (\textit{SHapley Additive exPlanations}) transpose aux modèles
prédictifs la théorie des jeux coopératifs : chaque variable est un joueur,
le score du modèle est le gain à partager. Pour une mesure donnée, la valeur
de Shapley de la variable $i$ est sa contribution marginale moyenne, moyennée
sur tous les sous-ensembles $S$ des autres variables :

\begin{equation}
\phi_i = \sum_{S \subseteq F \setminus \{i\}}
\frac{|S|!\,(|F|-|S|-1)!}{|F|!}\Big(f(S\cup\{i\}) - f(S)\Big)
\label{eq:shapley}
\end{equation}

où $F$ est l'ensemble des variables et $f(S)$ la prédiction obtenue en
n'utilisant que les variables de $S$. La propriété essentielle est
l'\textbf{additivité} : le score du modèle est la valeur de base $\phi_0$
(la prédiction moyenne du modèle) augmentée des contributions de toutes les
variables :

\begin{equation}
f(x) = \phi_0 + \sum_{i=1}^{n} \phi_i
\label{eq:additivite}
\end{equation}

Une valeur $\phi_i$ positive pousse la mesure vers « attaque », une valeur
négative vers « normal ». Le notebook d'explicabilité vérifie cette identité
sur une mesure isolée : la valeur de base (environ 0,5) augmentée des 53
contributions redonne exactement le score prédit par le modèle.

\subsection{Mise en œuvre dans le projet}

Le modèle retenu (Extra Trees, chapitre~\ref{chap4}) est un ensemble
d'arbres : nous utilisons \texttt{shap.TreeExplainer}, qui calcule les
valeurs de l'équation~\eqref{eq:shapley} de façon \emph{exacte} et rapide en
exploitant la structure des arbres, sans énumérer les sous-ensembles de
variables. Le calcul, intégré au script d'étude et au service de l'API,
suit toujours le même principe :

\begin{enumerate}
  \item \textbf{Lignes expliquées.} Seules des mesures falsifiées jamais
        vues à l'entraînement sont expliquées, hors zone de garde : on
        cherche comment le modèle \emph{reconnaît} une attaque, pas ce qu'il
        a appris par cœur. Pour l'étude de référence, ce sont les mesures
        d'attaque du test des cinq premières rotations (au plus 1\,000
        tirées par rotation, graine fixe), soit 4\,394 mesures au total ;
        après un réentraînement lancé depuis l'interface, le même calcul
        porte sur les attaques de la validation et du test (au plus 300).
  \item \textbf{Valeurs SHAP.} \texttt{TreeExplainer} produit une valeur
        par couple (mesure, variable).
  \item \textbf{Importance.} On prend la valeur absolue — le sens de
        l'effet importe peu, seule son intensité classe les variables —
        puis la moyenne sur les mesures expliquées :
        \begin{equation}
        I_j = \frac{1}{N}\sum_{k=1}^{N}\left|\phi_j^{(k)}\right|
        \label{eq:importance}
        \end{equation}
  \item \textbf{Classement.} Les variables sont triées par $I_j$
        décroissante : une fois sur toutes les mesures (importance globale,
        les 15 premières), une fois par scénario simulé en ne gardant que
        ses mesures (les 6 premières). Le regroupement par scénario utilise
        la colonne \texttt{attack\_type} du jeu de données, qui n'est jamais
        fournie au modèle.
\end{enumerate}

Le fichier \texttt{shap\_transmission\_reel.json} enregistre ces résultats ;
ils sont chargés dans la table \texttt{transmission\_shap} de la base de
données et affichés par la page « Explicabilité SHAP » du tableau de bord.
Le notebook \texttt{04\_shap.ipynb}, dans le répertoire
\texttt{notebook/transmission}, refait le calcul sur 500 mesures (graine
fixe) et vérifie que le classement est stable d'un tirage à l'autre.

\subsection{Importance globale des variables}

La figure~\ref{fig:expl_shap_global} et le
tableau~\ref{tab:expl_shap_global} présentent les dix variables les plus
importantes pour l'Extra Trees.

\begin{figure}[H]
\centering
\figtrois{fig_shap_global.png}{0.85}
\caption{Importance globale des variables (moyenne des valeurs SHAP
absolues), Extra Trees.}
\label{fig:expl_shap_global}
\end{figure}

\begin{table}[H]
\centering
\small
\renewcommand{\arraystretch}{1.3}
\setlength{\tabcolsep}{4pt}
\caption{Dix premières variables par importance SHAP globale (moyenne des
valeurs absolues).}
\label{tab:expl_shap_global}
\begin{tabularx}{\textwidth}{@{}|c|>{\ttfamily\footnotesize}l|r|>{\raggedright\arraybackslash}X|@{}}
\hline
\rowcolor{thead}
\textcolor{white}{\textbf{Rang}} & \textcolor{white}{\textbf{Variable}} &
\textcolor{white}{\textbf{SHAP}} &
\textcolor{white}{\textbf{Interprétation}} \\
\hline
1 & fridge\_repeat\_2min & 0,0609 & Séquence de 2 min du frigo déjà transmise
(rejeu) \\
\rowcolor{trowa}
2 & fridge\_max\_5min & 0,0540 & Maximum du frigo sur 5 min \\
3 & fridge\_q10\_5min & 0,0471 & Quantile 10\,\% du frigo sur 5 min \\
\rowcolor{trowa}
4 & fridge\_std\_1min & 0,0443 & Écart-type du frigo sur 1 min (gel,
injection) \\
5 & fridge\_mean\_5min & 0,0427 & Moyenne du frigo sur 5 min \\
\rowcolor{trowa}
6 & fridge\_mean\_15min & 0,0423 & Moyenne du frigo sur 15 min (biais,
dérive) \\
7 & fridge\_max\_15min & 0,0313 & Maximum du frigo sur 15 min \\
\rowcolor{trowa}
8 & fridge\_q90\_5min & 0,0233 & Quantile 90\,\% du frigo sur 5 min \\
9 & fridge\_q10\_15min & 0,0217 & Quantile 10\,\% du frigo sur 15 min \\
\rowcolor{trowa}
10 & fridge\_thermostat\_diff\_mean\_15min & 0,0216 & Écart frigo--thermostat
moyen sur 15 min \\
\hline
\end{tabularx}
\end{table}

Deux enseignements se dégagent. D'abord, toutes les variables de tête
concernent le \textbf{frigo} : c'est cohérent, puisque la majorité des neuf
scénarios falsifient cette sonde. Ensuite, la première place revient à la
répétition de séquence, conçue pour détecter le rejeu, et les suivantes
décrivent la \emph{forme} de la distribution des mesures sur 5 à 15 minutes
(maximum, quantiles, moyenne) : le modèle s'appuie sur des indices qui ont
un sens physique, pas sur des artefacts du jeu de données. Une précaution de
lecture : l'importance SHAP n'est pas une causalité, et deux variables
corrélées se partagent l'effet.

\subsection{Importance par scénario et cohérence physique}

Si le modèle a réellement appris la mécanique des attaques, chaque scénario
doit être reconnu par la variable conçue pour lui. Le
tableau~\ref{tab:expl_shap_scenario} donne, pour chaque scénario simulé, ses
trois variables les plus importantes ; la
figure~\ref{fig:expl_shap_scenario} montre le classement complet.

\begin{table}[H]
\centering
\small
\renewcommand{\arraystretch}{1.35}
\caption{Variables dominantes selon SHAP, par scénario d'attaque.}
\label{tab:expl_shap_scenario}
\begin{tabularx}{\textwidth}{@{}|>{\raggedright\arraybackslash}p{3.3cm}|>{\raggedright\arraybackslash}X|@{}}
\hline
\rowcolor{thead}
\textcolor{white}{\textbf{Scénario}} &
\textcolor{white}{\textbf{Trois premières variables (SHAP moyen)}} \\
\hline
Gel du frigo & écart-type 1 min (0,099) ; moyenne 5 min (0,077) ;
maximum 5 min (0,053) \\
\rowcolor{trowa}
Rejeu du frigo & répétition de séquence 2 min (\textbf{0,397}) ;
maximum 15 min (0,026) ; quantile 10\,\% du thermostat 15 min (0,022) \\
Injection nominale & écart-type 1 min (0,169) ; maximum 5 min (0,070) ;
écart-type 5 min (0,069) \\
\rowcolor{trowa}
Biais $-0{,}5$\,\textdegree C & maximum 5 min (0,048) ; quantile 10\,\% 5 min
(0,045) ; maximum 15 min (0,043) \\
Biais $-1$\,\textdegree C & moyenne 15 min (0,071) ; quantile 10\,\% 5 min
(0,065) ; maximum 5 min (0,064) \\
\rowcolor{trowa}
Biais $-2$\,\textdegree C & quantile 10\,\% 5 min (0,095) ; moyenne 5 min
(0,090) ; moyenne 15 min (0,081) \\
Dérive & quantile 10\,\% 5 min (0,072) ; maximum 5 min (0,064) ;
maximum 15 min (0,047) \\
\rowcolor{trowa}
Gel du thermostat & écart-type 1 min du thermostat (0,138) ; écart-type 5 min
(0,082) ; variation absolue 5 min (0,075) \\
\hline
\end{tabularx}
\end{table}

\begin{figure}[H]
\centering
\figtrois{fig_shap_scenario.png}{0.95}
\caption{Importance des variables par scénario d'attaque.}
\label{fig:expl_shap_scenario}
\end{figure}

La lecture du tableau~\ref{tab:expl_shap_scenario} confirme la cohérence
physique du modèle :

\begin{itemize}
  \item le \textbf{rejeu} est reconnu presque uniquement par la répétition de
        séquence (0,397, environ quinze fois plus que la variable suivante) ;
  \item le \textbf{gel} et l'\textbf{injection nominale} sont reconnus par
        l'écart-type à 1 minute, qui s'effondre lorsque la mesure est figée
        ou trop régulière ; pour le gel du thermostat, ce sont les
        écarts-types du thermostat qui dominent ;
  \item les \textbf{biais} et la \textbf{dérive} reposent sur la forme de la
        distribution sur 5 à 15 minutes (quantile 10\,\%, maximum, moyennes) :
        un décalage de quelques dixièmes de degré déplace le haut et le bas
        de la distribution bornée du frigo (1--14\,\textdegree C) avant que
        la moyenne, noyée dans un bruit de $\sigma \approx
        3{,}4$\,\textdegree C par mesure, ne bouge.
\end{itemize}

Pour les petits biais ($-0{,}5$ et $-1$\,\textdegree C) et la dérive, les
importances sont plus faibles et étalées entre plusieurs variables : le
modèle ne dispose pas d'un indice net. Cette observation annonce les
résultats de la section suivante, où ces scénarios sont les moins bien
détectés. Elle montre aussi que les indices ne séparent pas un biais d'une
dérive, point repris à la section~\ref{sec:expl_deduction}.

\subsection{Explication locale d'une alerte}

SHAP explique aussi \emph{une} alerte précise : pour une mesure donnée, les
contributions positives (vers « attaque ») et négatives (vers « normal »)
montrent quelles variables ont fait passer le score au-dessus du seuil
d'alerte. La figure~\ref{fig:expl_shap_local} prend une mesure de gel :
l'opérateur y lit que l'alerte vient surtout de l'effondrement de l'écart-type
à 1 minute. C'est cette lecture, reprise par scénario dans le rapport, qui
justifie chaque alerte.

\begin{figure}[H]
\centering
\figtrois{fig_shap_local.png}{0.90}
\caption{Explication locale d'une alerte de gel : contributions des variables
au score de la mesure.}
\label{fig:expl_shap_local}
\end{figure}


% =====================================================================
\section{Rapport de recommandations pour l'opérateur}
\label{sec:expl_rapport}

\subsection{Objectif et génération}

Le rapport transforme les résultats du modèle en \textbf{consignes lisibles
par un non-spécialiste}. Il est produit par la fonction
\texttt{operator\_report} du service de transmission, servie par la route
\texttt{GET /api/v1/transmission/report} et affiché par la page « Rapport
pour l'opérateur » du tableau de bord. Pour chacun des neuf scénarios
simulés, il répond à cinq questions : que fait l'attaquant, le modèle réagit-il
bien, sur quels indices s'appuie-t-il, quelle famille d'attaque ces indices
laissent-ils deviner, et que doit faire l'opérateur. Le modèle étant
binaire, le scénario sert ici de regroupement des résultats d'évaluation, et
non de réponse du modèle.

Aucun chiffre n'est écrit en dur : tous sont lus dans les résultats
enregistrés — l'étude de référence (20 rotations) par défaut, ou le dernier
entraînement lancé depuis l'interface — ce qui garantit que le rapport
reste cohérent après chaque réentraînement. Il est produit en deux formats,
un document JSON structuré et son rendu Markdown téléchargeable, en suivant
les étapes suivantes :

\begin{enumerate}
  \item \textbf{choix du modèle en service} : le meilleur du classement sur
        la \emph{validation} (jamais sur le test), à l'exclusion de l'ancien
        modèle de référence ; c'est l'Extra Trees (chapitre~\ref{chap4}) ;
  \item \textbf{description du réglage d'alerte} retenu sur la validation :
        un seuil de score et une règle de persistance (alerte lorsque au
        moins $k$ mesures sur les $n$ dernières dépassent le seuil) ;
  \item \textbf{calcul, pour chaque scénario,} de la part des épisodes
        détectés, de la part des mesures falsifiées signalées et du délai
        médian de détection, mesurés au test ;
  \item \textbf{attribution d'un niveau de détection}, selon les critères du
        tableau~\ref{tab:expl_niveaux} ;
  \item \textbf{ajout des indices du modèle} : les trois premières variables
        SHAP du scénario, traduites en français lisible (par exemple
        « écart-type sur 1 min du frigo ») ;
  \item \textbf{déduction de la famille d'attaque} à partir de ces seules
        variables (section~\ref{sec:expl_deduction}), comparée à la famille
        réellement simulée, puis ajout des consignes de la famille ;
  \item \textbf{ajout des limites connues et des règles générales}, puis
        assemblage du document.
\end{enumerate}

\subsection{Le modèle en service}

Le rapport commence par présenter le modèle retenu. Le
tableau~\ref{tab:expl_modele} reprend ses performances au test, cumulées sur
les 20 rotations de l'étude.

\begin{table}[H]
\centering
\small
\renewcommand{\arraystretch}{1.3}
\caption{Modèle en service et performances au test (20 rotations).}
\label{tab:expl_modele}
\begin{tabular}{|l|r|}
\hline
\rowcolor{thead}
\textcolor{white}{\textbf{Indicateur}} & \textcolor{white}{\textbf{Valeur}} \\
\hline
Modèle & Extra Trees \\
\rowcolor{trowa}
Réglage d'alerte & seuil 0,70 ; 3 mesures sur 6 (30 s) \\
Épisodes détectés & 70 sur 73 \\
\rowcolor{trowa}
Rappel (mesures falsifiées signalées) & 65,2\,\% \\
Précision & 93,9\,\% \\
\rowcolor{trowa}
Score F1 & 0,77 \\
Fausses alertes par heure & 0,43 \\
\rowcolor{trowa}
Délai médian de détection & 3,9 min \\
\hline
\end{tabular}
\end{table}

Le rapport traduit ces chiffres en une phrase accessible : sur 100 mesures
signalées, environ 94 sont réellement falsifiées, et environ 65\,\% des
mesures falsifiées sont signalées. Le réglage respecte les contraintes fixées
sur la validation — précision d'au moins 90\,\% et au plus 0,5 fausse alerte
par heure — afin de ne pas noyer l'opérateur sous les fausses alertes
(chapitre~\ref{chap4}).

\subsection{Niveau de détection par scénario}

Un épisode compte comme détecté dès qu'une alerte tombe pendant sa durée.
Cette mesure seule serait trop indulgente : un scénario peut avoir tous ses
épisodes détectés tout en ne signalant qu'une petite part des mesures
falsifiées, ou trop tard. Le niveau combine donc trois mesures, selon le
tableau~\ref{tab:expl_niveaux}. Ces seuils sont des critères d'acceptation
fixés a priori pour rendre les résultats lisibles ; ils ne sont pas calculés
à partir des données et ne servent à aucun choix du modèle.

\begin{table}[H]
\centering
\small
\renewcommand{\arraystretch}{1.35}
\caption{Critères du niveau de détection d'un scénario (mesurés au test).}
\label{tab:expl_niveaux}
\begin{tabularx}{\textwidth}{@{}|>{\bfseries}p{4.2cm}|>{\raggedright\arraybackslash}X|@{}}
\hline
\rowcolor{thead}
\textcolor{white}{\textbf{Niveau}} & \textcolor{white}{\textbf{Condition}} \\
\hline
Bien détecté & Au moins 85\,\% des épisodes détectés, au moins 70\,\% des
mesures falsifiées signalées et délai médian d'au plus 5 min. \\
\hline
\rowcolor{trowa}
Détecté partiellement & Moins de 85\,\% des épisodes détectés, ou moins de
70\,\% des mesures falsifiées signalées. \\
\hline
Détecté tardivement & Délai médian supérieur à 5 min. \\
\hline
\rowcolor{trowa}
Mal détecté & Moins de 60\,\% des épisodes détectés : vérification humaine
indispensable. \\
\hline
\end{tabularx}
\end{table}

Le tableau~\ref{tab:expl_resultats_scenarios} applique ces critères aux
neuf scénarios.

\begin{table}[H]
\centering
\small
\renewcommand{\arraystretch}{1.3}
\caption{Niveau de détection de chaque scénario au test (Extra Trees).}
\label{tab:expl_resultats_scenarios}
\begin{tabularx}{\textwidth}{@{}|>{\raggedright\arraybackslash}X|c|c|c|>{\raggedright\arraybackslash}p{3.3cm}|@{}}
\hline
\rowcolor{thead}
\textcolor{white}{\textbf{Scénario}} & \textcolor{white}{\textbf{Épisodes}} &
\textcolor{white}{\textbf{Mesures}} & \textcolor{white}{\textbf{Délai}} &
\textcolor{white}{\textbf{Niveau}} \\
\hline
Gel du frigo & 8/8 & 95,0\,\% & 1,0 min & bien détecté \\
\rowcolor{trowa}
Rejeu du frigo & 8/9 & 73,8\,\% & 2,1 min & bien détecté \\
Injection nominale & 9/9 & 88,7\,\% & 2,2 min & bien détecté \\
\rowcolor{trowa}
Biais $-0{,}5$\,\textdegree C & 9/10 & 26,9\,\% & 11,0 min & partiel et
tardif \\
Biais $-1$\,\textdegree C & 8/8 & 52,3\,\% & 6,0 min & partiel et tardif \\
\rowcolor{trowa}
Biais $-2$\,\textdegree C & 7/7 & 83,9\,\% & 3,0 min & bien détecté \\
Dérive & 7/7 & 31,8\,\% & 11,6 min & partiel et tardif \\
\rowcolor{trowa}
Gel du thermostat & 7/7 & 83,4\,\% & 3,6 min & bien détecté \\
Biais du thermostat & 7/8 & 51,8\,\% & 6,9 min & partiel et tardif \\
\hline
\end{tabularx}
\end{table}

Cinq scénarios sur neuf sont bien détectés : le gel, le rejeu, l'injection
nominale, le biais de $-2$\,\textdegree C et le gel du thermostat. Les petits
biais, la dérive et le biais du thermostat sont détectés partiellement et
tardivement : le décalage est inférieur au bruit naturel du capteur, et seule
la forme de la distribution sur 5 à 15 minutes trahit l'attaque. Le rapport
signale explicitement cette limite à l'opérateur, qui ne doit pas se reposer
sur le seul détecteur pour ces cas.

\subsection{Consignes par famille d'attaque}

Les consignes dépendent du \emph{mécanisme} de l'attaque, pas du capteur
visé : le même geste convient face à un gel du frigo ou du thermostat. Elles
sont rédigées une fois par famille ; deux scénarios d'une même famille, par
exemple les trois biais du frigo, affichent donc les mêmes actions alors que
leur niveau de détection diffère. Le tableau~\ref{tab:expl_consignes} résume
les cinq familles.

\begin{table}[H]
\centering
\small
\renewcommand{\arraystretch}{1.35}
\caption{Consignes de l'opérateur par famille d'attaque.}
\label{tab:expl_consignes}
\begin{tabularx}{\textwidth}{@{}|>{\raggedright\arraybackslash}p{2.7cm}|>{\raggedright\arraybackslash}X|>{\raggedright\arraybackslash}X|@{}}
\hline
\rowcolor{thead}
\textcolor{white}{\textbf{Famille}} & \textcolor{white}{\textbf{Risque}} &
\textcolor{white}{\textbf{Actions principales}} \\
\hline
Gel (mesure figée) & La même valeur est reçue en boucle : une vraie dérive
passerait inaperçue. & Relever la température sur place ; vérifier la liaison
capteur--passerelle--broker et redémarrer le capteur ; isoler le lot si la
mesure réelle est hors 2--8\,\textdegree C. \\
\hline
\rowcolor{trowa}
Rejeu & Des mesures passées sont renvoyées : l'état réel du frigo est
masqué. & Relever la température sur place ; vérifier les horodatages des
messages MQTT ; activer TLS et l'authentification sur le broker ; isoler le
capteur. \\
\hline
Injection nominale & Des valeurs forgées proches de la consigne remplacent
les vraies mesures : le frigo paraît parfait quoi qu'il arrive. & Relever la
température sur place ; vérifier l'intégrité du firmware ; activer TLS et
l'authentification ; bloquer la libération des lots (BPF). \\
\hline
\rowcolor{trowa}
Biais & Le frigo paraît plus froid qu'il ne l'est : le produit peut sortir
de la plage 2--8\,\textdegree C sans alerte. & Comparer avec un thermomètre
de référence étalonné ; recalibrer ou remplacer le capteur si l'écart est
confirmé ; vérifier la passerelle. \\
\hline
Dérive & Le décalage augmente peu à peu et finit par cacher une sortie de
plage. & Comparer avec un thermomètre de référence ; examiner l'historique de
la journée ; recalibrer le capteur ; bloquer la libération des lots (BPF). \\
\hline
\end{tabularx}
\end{table}

Toutes les familles commencent par la même action : \textbf{vérifier sur
place avec un thermomètre de contrôle}. Le détecteur signale en effet une
transmission suspecte, pas une température.

\subsection{Déduire la famille d'attaque à partir des indices}
\label{sec:expl_deduction}

Le modèle ne fournit qu'un score d'attaque : il ne dit pas « ceci est un
gel ». Pour associer une consigne à une alerte, la fonction
\texttt{infer\_families} parcourt les variables SHAP dominantes, de la plus
forte à la plus faible, et applique des règles fondées sur la nature de la
variable (tableau~\ref{tab:expl_deduction}).

\begin{table}[H]
\centering
\small
\renewcommand{\arraystretch}{1.35}
\caption{Règles de déduction de la famille d'attaque à partir des variables
dominantes.}
\label{tab:expl_deduction}
\begin{tabularx}{\textwidth}{@{}|>{\raggedright\arraybackslash}X|>{\raggedright\arraybackslash}X|@{}}
\hline
\rowcolor{thead}
\textcolor{white}{\textbf{Indice dominant}} &
\textcolor{white}{\textbf{Famille(s) proposée(s)}} \\
\hline
Séquence déjà transmise parmi les trois premières variables & Rejeu \\
\hline
\rowcolor{trowa}
Dispersion en tête (écart-type, variation d'une mesure à l'autre) & Gel ou
injection nominale. L'injection nominale vient en premier si l'écart-type sur
5 min ou la variation absolue figure parmi les quatre premières variables,
sinon le gel. \\
\hline
Niveau en tête (moyenne, maximum, quantiles) & Biais ou dérive \\
\hline
\rowcolor{trowa}
Aucune règle applicable & Famille non identifiée : consigne générale \\
\hline
\end{tabularx}
\end{table}

Deux précautions s'imposent. D'abord, les indices ne séparent pas toujours
deux familles voisines : un gel et une injection nominale ont tous deux un
écart-type à 1 minute très bas, et un biais et une dérive déplacent les mêmes
quantiles et les mêmes moyennes. La déduction renvoie alors \emph{deux}
hypothèses, la plus probable en premier. Ensuite, il s'agit d'une hypothèse,
jamais d'un diagnostic : la fiche porte l'avertissement « à confirmer par
une vérification sur place », et la consigne reste dans tous les cas de
relever la température avec un thermomètre de contrôle.

Le rapport teste lui-même cette déduction. Pour chaque scénario simulé, il
applique les règles aux seules variables SHAP du scénario et vérifie si la
famille réellement simulée figure parmi les hypothèses
(tableau~\ref{tab:expl_deduction_resultats}).

\begin{table}[H]
\centering
\small
\renewcommand{\arraystretch}{1.3}
\caption{Famille déduite des indices SHAP, comparée à la famille simulée.}
\label{tab:expl_deduction_resultats}
\begin{tabularx}{\textwidth}{@{}|>{\raggedright\arraybackslash}X|>{\raggedright\arraybackslash}X|c|@{}}
\hline
\rowcolor{thead}
\textcolor{white}{\textbf{Scénario simulé}} &
\textcolor{white}{\textbf{Hypothèses proposées}} &
\textcolor{white}{\textbf{Famille simulée retrouvée}} \\
\hline
Gel du frigo & gel ; injection nominale & oui (1\textsuperscript{re}) \\
\rowcolor{trowa}
Rejeu du frigo & rejeu & oui (1\textsuperscript{re}) \\
Injection nominale & injection nominale ; gel & oui (1\textsuperscript{re}) \\
\rowcolor{trowa}
Biais $-0{,}5$\,\textdegree C & biais ; dérive & oui (1\textsuperscript{re}) \\
Biais $-1$\,\textdegree C & biais ; dérive & oui (1\textsuperscript{re}) \\
\rowcolor{trowa}
Biais $-2$\,\textdegree C & biais ; dérive & oui (1\textsuperscript{re}) \\
Dérive & biais ; dérive & oui (2\textsuperscript{e}) \\
\rowcolor{trowa}
Gel du thermostat & injection nominale ; gel & oui (2\textsuperscript{e}) \\
Biais du thermostat & --- & non évalué \\
\hline
\end{tabularx}
\end{table}

La bonne famille figure parmi les hypothèses pour les huit scénarios
évalués, et en première position pour six d'entre eux. Le biais du
thermostat n'est pas évalué faute de valeurs SHAP enregistrées pour ce
scénario. Ce résultat est mesuré sur des attaques simulées dont les
variables dominantes sont connues ; il montre que les indices portent
l'information nécessaire, sans garantir la même réussite sur une attaque
réelle inédite. La perspective naturelle est un modèle multiclasse, entraîné
à prédire aussi la famille, au prix d'un nouvel entraînement et d'une
nouvelle évaluation.

\subsection{Exemple de fiche et page du tableau de bord}

L'encadré suivant présente la fiche générée pour le scénario de gel du
frigo : les indices du modèle proviennent de l'analyse SHAP et les chiffres
des résultats au test.

\begin{center}
\fbox{\parbox{0.93\textwidth}{\small
\textbf{Mesure figée (MITM) --- MITM\_gel (frigo)}\\[2pt]
\textbf{Ce que fait l'attaquant :} la dernière mesure du frigo est répétée en
boucle.\\
\textbf{Détection :} épisodes 8/8, délai médian 1,0 min, 95\,\% des mesures
signalées --- \textbf{bien détecté}.\\
\textbf{Ce que le modèle regarde :} écart-type sur 1 min du frigo ; moyenne
sur 5 min du frigo ; maximum sur 5 min du frigo.\\
\textbf{Famille déduite de ces indices :} mesure figée (MITM) ou injection de
valeurs plausibles --- famille simulée retrouvée.\\
\textbf{Risque :} le serveur reçoit la même valeur en boucle ; une vraie
dérive de température passerait inaperçue.\\
\textbf{À faire :} (1) relever la température sur place avec un thermomètre
de contrôle ; (2) vérifier la liaison capteur, passerelle et broker, puis
redémarrer le capteur ; (3) si la mesure réelle est hors
2--8\,\textdegree C, isoler le lot et ouvrir un incident chaîne du froid ;
(4) comparer l'horodatage des messages MQTT avec l'horloge du capteur.}}
\end{center}

La figure~\ref{fig:expl_capture_rapport} montre la page « Rapport pour
l'opérateur » du tableau de bord, qui affiche le modèle en service, le bloc
rappelant que le modèle ne nomme pas l'attaque, le critère du niveau, les
fiches de chaque scénario avec la famille déduite, les limites et les règles
générales.

\begin{figure}[H]
\centering
\figtrois{rapport_recommandations.png}{0.95}
\caption{Page du rapport de recommandations dans le tableau de bord.}
\label{fig:expl_capture_rapport}
\end{figure}

\subsection{Recommandation d'une alerte précise}

Pour une alerte en temps réel, la fonction \texttt{alert\_recommendation}
(route \texttt{POST /transmission/report}) produit une consigne ciblée à
partir de la probabilité de l'alerte, des variables SHAP dominantes et,
facultativement, du capteur concerné. Le modèle étant binaire, le type
d'attaque n'est pas une entrée obligatoire :

\begin{itemize}
  \item sans type indiqué, la famille est \emph{déduite} des variables
        dominantes (section~\ref{sec:expl_deduction}) ; la fiche porte
        l'avertissement « hypothèse à confirmer sur place », la mention
        « famille probable » et les autres hypothèses éventuelles ; le
        capteur est lu dans le préfixe de la première variable
        (\texttt{fridge\_}, \texttt{thermostat\_}) ;
  \item si le type est indiqué (scénario connu, par exemple en simulation), il
        est repris tel quel ;
  \item si aucune règle ne s'applique, la fiche annonce une « attaque de
        transmission suspectée (type non identifié) » et donne les règles
        générales de la section suivante.
\end{itemize}

La fonction y ajoute un niveau de \textbf{confiance} déduit de la
probabilité (tableau~\ref{tab:expl_confiance}), puis la consigne : contrôler
la température sur place avant toute décision, puis appliquer l'action
adaptée à la famille retenue.

\begin{table}[H]
\centering
\small
\renewcommand{\arraystretch}{1.3}
\caption{Niveau de confiance associé à la probabilité d'une alerte.}
\label{tab:expl_confiance}
\begin{tabular}{|l|l|}
\hline
\rowcolor{thead}
\textcolor{white}{\textbf{Probabilité}} & \textcolor{white}{\textbf{Confiance}} \\
\hline
supérieure ou égale à 0,80 & élevée \\
\rowcolor{trowa}
entre 0,50 et 0,80 & moyenne \\
inférieure à 0,50 & faible \\
\hline
\end{tabular}
\end{table}

\subsection{Limites connues et règles générales}

Le rapport se termine par les limites du détecteur, écrites pour
l'opérateur. La première concerne les petits biais du frigo ($-0{,}5$ et
$-1$\,\textdegree C) : le décalage est plus petit que le bruit du capteur
($\sigma \approx 3{,}4$\,\textdegree C par mesure), la détection vient de la
forme de la distribution sur 5 à 15 min et reste plus lente et moins sûre
que pour les autres scénarios. L'étude complémentaire sur les biais longs
(épisodes de 1 à 3 h) le confirme : avec les mêmes variables, les biais de
$-0{,}5$ et $-1$\,\textdegree C sont détectés dans 15 épisodes sur 15
chacun, mais avec un délai médian de 8,1 et 9,1 min et un seuil plus
permissif (1,47 fausse alerte par heure au test) ; ajouter des moyennes
glissantes sur 30 et 60 minutes n'améliore rien.

La seconde limite tient à la nature binaire du modèle : la famille d'attaque
n'est qu'une hypothèse déduite des indices, que la vérification sur place
confirme ou corrige.

Les règles de conduite valables pour toutes les attaques sont ensuite
rappelées :

\begin{itemize}
  \item une alerte se vérifie d'abord sur place, avec un thermomètre de
        contrôle ;
  \item un lot ne se libère jamais sur la seule base d'une mesure
        contestée ;
  \item le capteur est isolé du réseau et le broker vérifié (TLS,
        authentification, horodatage) dès que l'attaque est confirmée ;
  \item chaque alerte, la vérification faite et la décision prise sont
        journalisées.
\end{itemize}


% =====================================================================
\section{Conclusion}

Ce chapitre a répondu aux deux questions de l'opérateur. L'analyse SHAP,
calculée de façon exacte sur le modèle à arbres, montre que l'Extra Trees
s'appuie sur des indices cohérents avec la mécanique des attaques : la
répétition de séquence pour le rejeu, l'écart-type à 1 minute pour le gel
et l'injection, la forme de la distribution (quantiles, maximum, moyennes)
pour les biais et la dérive. Le rapport de recommandations utilise ces
indices, avec les résultats mesurés au test, pour attribuer à chaque
scénario simulé un niveau de détection et une consigne.

Le modèle étant binaire, il ne reconnaît pas le type d'attaque : la famille
est déduite des variables dominantes, et la bonne famille figure parmi les
hypothèses pour les huit scénarios évalués (six fois en première position).
Cette déduction reste une hypothèse, que la vérification sur place confirme
ou corrige ; entraîner un modèle multiclasse est la perspective qui
supprimerait cette limite.

Le bilan est nuancé et assumé : cinq scénarios sur neuf sont bien détectés,
tandis que les petits biais, la dérive et le biais du thermostat restent
détectés partiellement et tardivement. Le rapport l'indique à l'opérateur,
qui conserve la vérification sur place comme premier réflexe. Le chapitre
suivant décrit la réalisation de l'API et du tableau de bord qui exposent
ces résultats.
